\documentclass[../unravbook.tex]{subfiles}

\usepackage{parskip}
\setlength{\parskip}{10pt}
\setstretch{1.15}

\setenotez{
  list-heading = {\chapter*{#1}\markboth{#1}{#1}}
}

\begin{document}

\setcounter{chapter}{3}
\setcounter{endnote}{0}

\chapter{Dehumanized} %% {{{
    %% end header ----------------------------
    


Since the \enquote{Al--Aqsa Flood} operation started on October 7, 2023, an informed consensus has emerged that apart from whatever \enquote{paradigm shift in the philosophy of the resistance} may have been envisioned,\footnote{Muhammed Hüseyin Mercan, \enquote{Operation al-Aqsa Flood: A Rupture in the History of the Palestinian Resistance and Its Implications,} January 1, 2024, SETA Foundation, \url{https://www.setav.org/en/operation-al-aqsa-flood-a-rupture-in-the-history-of-the-palestinian-resistance-and-its-implications}.}  its most consequential outcome over the long term has been increased outrage in Western societies toward Israel and sympathy --- or even outright support --- for its regional adversaries.\footnote{\enquote{Net Favourability towards Israel Reaches New Lows in Key Western European Countries.} YouGov, June 3, 2025, \url{https://yougov.com/en-gb/articles/52279-net-favourability-towards-israel-reaches-new-lows-in-key-western-european-countries}.} This ostensibly counter-intuitive result has been attributed to multiple factors, such as the view that the operation was either facilitated by Israeli authorities or the result of their own actions;\footnote{Jonathan Cook, \enquote{Why Is Western Media Ignoring Evidence of Israel's Own Actions on 7 October?,} Middle East Eye, December 14, 2023. \url{https://www.middleeasteye.net/opinion/israel-palestine-war-media-ignoring-evidence-actions-7-october.}.} or that it is unsurprising given decades of Israeli oppression;\footnote{Diego, García-Sayan \enquote{International Justice and Genocide in Gaza.} Ediciones EL PAÍS S.L., January 19, 2024 \url{https://english.elpais.com/opinion/2024-01-19/international-justice-and-genocide-in-gaza.html}.} or that Israel's targeting afterward of vulnerable civilians and its vast destruction of Gaza's infrastructure cannot be overlooked, as it explains why Hamas took such extreme measures in the first place.\footnote{NATO Strategic Communications Centre of Excellence Riga (StratCom COE), 2015, Hamas' use of human shields in Gaza, 2008--2014. \url{https://stratcomcoe.org/cuploads/pfiles/hamas_human_shields.pdf}. Similar tactics were used by the FLN in Algeria in the 1950s-60s. Lincoln Krause, \enquote{\enquote*{A War to the Death}: The Ugly Underside of an Iconic Insurgency.} War on the Rocks, April 19, 2019. \url{https://warontherocks.com/a-war-to-the-death-the-ugly-underside-of-an-iconic-insurgency/}.} This shift could also reflect longer-term antecedents, such as the effects of algorithmic-driven social media on young people's attitudes\footnote{Lukas Hess and Sarah Häusermann, \enquote{The Quest for Exposure: How pro-Palestine Activists Navigate Algorithmic Visibility and Content Moderation on TikTok}, Platforms \& Society 3 (March 2026). \url{https://doi.org/10.1177/29768624261427203}.} and their increasing receptivity to any social justice ideological framework that one apologist for the operation described as \enquote{de-exceptionaliz[ing] Palestine by harnessing the Palestinian cause's universalist appeal.}\footnote{Zahi Zalloua, \enquote{Against Exceptionalism,} Humanities 13, 2, 2024 \url{https://doi.org/10.3390/h13020050}.} 

Regardless, this outcome has presented jihadist MENA state and non-state actors with an opportunity to expand their doctrine of ``perpetual resistance'' into the social activism throughout Western Europe and North America.\footnote{Ramzy Baroud, \enquote{Armed Vs. Peaceful Resistance --- What You Need to Know about Muqawama in Gaza,} June 25, 2024, Counterpunch. \url{https://www.counterpunch.org/2024/06/25/armed-vs-peaceful-resistance-what-you-need-to-know-about-muqawama-in-gaza/}.} They were never shy before about heaping Israel with scorn and shame in the MENA region but now, neither were those supporting them in those wider geographies. A sophisticated follow-up campaign to discredit Israel's international standing quickly engulfed Western milieus.\footnote{E.g., Brett Stephens, The Year American Jews Woke Up, Oct. 4, 2024. New York Times. \url{https://www.nytimes.com/2024/10/04/opinion/israel-jews-antisemitism.html}.} Its much higher tempo after October 7, particularly on college campuses and in youth-oriented media, naturally generated considerable alarm.\footnote{See, e.g., the theatrical file \textit{October 8}. \enquote{October 8} | Official Website, March 14, 2025. \url{https://www.october8film.com/}.}  While silent about the recent atrocities committed by Israel's enemies, the incessant fomenting of outrage toward Israelis and anyone who supported them --- derisively labeled as \enquote{Zionists} --- gave strong indications of being part of a cohesive, overall strategy. Claims of an \enquote{impending genocide}\footnote{E.g., Center for Constitutional Rights. \enquote{Israel's Unfolding Crime of Genocide of the Palestinian People \& U.S. Failure to Prevent and Complicity in Genocide,} October 18, 2023. \url{https://ccrjustice.org/israel-s-unfolding-crime-genocide-palestinian-people-us-failure-prevent-and-complicity-genocide}; Sai Englert, \enquote{Impending Genocide,} Sidecar, October 16, 2023, \url{https://newleftreview.org/sidecar/posts/impending-genocide}. (``In Gaza, Israel is gearing up to commit genocide.'').} filled the Infosphere as \enquote{within minutes, graphic imagery and bombastic propaganda began to flood social media platforms,}\footnote{P. W. Singer and Emerson T. Brooking, {Gaza and the Future of Information Warfare,} Foreign Affairs, November 15, 2023. \url{https://www.foreignaffairs.com/israel/gaza-and-future-information-warfare}.} well before any response from Israel actually took place. Claims of ever more disproportionate brutality and oppression on Israel's part was a recurring feature of the messaging\footnote{The Dahiya Doctrine \& Israel's Use of Disproportionate Force. July 31, 2024. IMEU. \url{https://imeu.org/resources/resources/explainer-the-dahiya-doctrine-israels-use-of-disproportionate-force/175}.} as NGOs and activist groups avidly predicted Israeli war crimes.\footnote{Lexis Lonas Cochran, \enquote{Students for Justice in Palestine becomes flashpoint in college antisemitism, free speech debate,} The Hill, November 16, 2023. \url{https://thehill.com/homenews/education/4311746-students-justice-palestine-antisemitism-free-speech-hamas-israel/} (\enquote{SJP wasted no time after the Oct. 7 Hamas terrorist attack on Israel to spring into action, organizing a \enquote*{Day of Resistance}. Chapters of the group said the world had witnessed \enquote*{a historic win for the Palestinian resistance.}})} Then, once Israel began responding, unconditional uniformly pleading, to use the phrasing of one such activist group, that the World \enquote{take moral action and implement an immediate ceasefire now.}\footnote{If Not Now (@ifNotNowOrg), November 13, 2023, X.com, \url{https://x.com/IfNotNowOrg/status/1724237701631656201?s=20}.}

The purpose of this vast onslaught of the anti-Israel pen in addition to the sword was never obscure. But, while many in Western societies could see that Israelis suffered immensely on October 7, with even greater physical and psychic trauma than what Americans experienced on September 11, 2001 and Europeans on November 13, 2015, they may not have realized they were now targets in a war of their own. Its was choreographed by either careless or deliberately misleading news reporting, along with strident \enquote{commentary} about what was actually going on in Gaza and the north of Israel, where Hezbollah \enquote{launched guided rockets and artillery onto three posts in the Shebaa Farms `in solidarity' with the Palestinian people.}\footnote{  Reuters, {Israel, Hezbollah Exchange Artillery, Rocket Fire,} October 8, 2023. \url{https://www.reuters.com/world/middle-east/israel-strikes-lebanon-after-hezbollah-hits-shebaa-farms-2023-10-08/}.} Eventually, and perhaps inevitably, the ability of major news outlets to stop mass incitement diminished, replaced by such things as covering with comment things such as a musical chant at a massive UK concert of \enquote{Death, death, death to the IDF.}\footnote{Annabel Sinclair, \enquote{BBC Says Bob Vylan's Glastonbury Set 'Utterly Unacceptable' as It Promises Review of Live Event Rules.} Jewish News, June 30, 2025. \url{https://www.jewishnews.co.uk/bbc-says-bob-vylans-glastonbury-set-utterly-unacceptable-as-it-promises-review-of-live-event-rules/}.} Provocations exploded across scheduled public events, campus encampments and supposedly spontaneous activist demonstrations and \enquote{actions,} often crossing over (or at least bordering) what might be legally constrained in the interest of public safety and civil order. All such tactics were intended to arouse the interest of people who might otherwise prefer to remain largely oblivious to the overall situation, moving them away from vague feelings of sympathy for those killed or taken hostage and toward hardened attitudes that compelled them to join in what was constantly depicted as the World's most important, righteous and still unresolved humanitarian cause: Palestine.\footnote{Noura Erekat, \enquote{Reflections from the Present on the Future of Political Action for Palestine,} MERIP, January 11, 2024, \url{https://www.merip.org/2024/01/four-reflections-on-political-action-and-palestine/}. (\enquote{Israel faces charges of genocide at the International Court of Justice, the International Criminal Court and in the eyes of millions marching in the streets of Tokyo, Seoul, London, Dublin, Cairo, Amman, Sydney, Chicago, San Francisco, New York, Toronto and other cities around the world. It is an unprecedented demonstration of robust solidarity with Palestinian liberation and pointed critiques of Israel that far exceed what existed before October 7.}} 

Unsurprisingly, the coordinated histrionics of this overall turn of events have had a profound impact on popular sentiment. Emblematic of its changed shape might be the social media views by two otherwise random mid-career professionals. A non-profit communications manager declared in early 2024 that after \enquote{11 years of demanding a Free Palestine \dots [a]partheid has led to genocide, and as Americans, we are complicit.}\footnote{Jaime Puente, February 27, 2024. Instagram post. \url{https://www.instagram.com/p/C32teUCOtU6/}.} A philosophy professor focused on the \enquote{nature of empathy}\footnote{Amy Coplan, Cal State Fullerton (n.d.). \url{https://philosophy.fullerton.edu/people/}.} shared in 2026 that while she cared ``deeply'' about \enquote{the people of Israel} and loathed Hamas, \enquote{the actions of a genocidal Israeli administration} were \enquote{unconscionable.}\footnote{Amy Coplan (@amycoplan), X, June 9, 2026, \url{https://x.com/amycoplan/status/2064585824445104286}; June 10, 2026, \url{https://x.com/amycoplan/status/2064773323406659865}.} Across their generation, a view has gradually taken hold that Israel thrives at the expense of Palestinian suffering because its \enquote{far-right} leaders (if not its ordinary citizens) can't resist doing genocidal things. This takes precedence no matter how impressive the nation's advances in the fields of science, technology or business may be, or how important it may be to Jewish communities worldwide. Thus, a certain habit of mind can now be seen as having taken root in the contemporary Western imaginary, inculcated by the litany of attacks on Israel, so that is now seen as a one-dimensional comic book villain, imposing its evil over the erased, \enquote{indigenous} topography of Palestine.

This perceptual reality encapsulates how accusations of Israeli apartheid have served as a predicate for further claims of Israeli genocide.\footnote{It is worth noting that claims of Israel's genocidal nature did not begin until after October 7. \parencite[59, (\enquote{I consider the Nakba itself to be an act of genocide.})]{Abdo-2018} However, the extent of such claims has exploded since October 7 and is even extended to countries in the West that express a willingness to tolerate Israel's continued existence \parencite[ (\enquote{As it has happened for the past five centuries, the Palestinian Genocide shows that European settler colonialism continues to perform one of its historical tasks: genocide. Empires and former empires like the US, the UK, and France continue to be authors and accomplices of genocide by enabling the genocidal aggression of a people.}]{Barreto-2024}.}  The two charges are worthy of further scrutiny together because essentially, they amount to a concatenated pair. Feeding off each other in their irrebuttable depiction of horrific intent, they are now the dual touchstones in attempts to dehumanize Israel as the perpetrator of crimes against humanity, leaving only misery in her wake. Thus, apart from labeling people (mainly Jews) as \enquote{apartheid loving Zionists}\footnote{\enquote{Appropriation of Rage: Why Allies/Accomplices Need to Stay in their Lane}, Medium, September 1, 2020. \url{https://medium.com/@avirden86/appropriation-of-rage-why-allies-accomplices-need-to-stay-in-their-lane-489b39c591c}.} and governments and organizational entities that do not oppose Israel's ongoing existence (mainly in Western socieities)  as ``accomplices of genocide,'' the hundreds of thousands of academics, media influencers, public figures and \enquote{just plain folks} who have chosen to utter these pronouncements enable the success of what amounts to the \enquote{master narrative} of Palestine. In its overarching essence, the collective force of such messaging is nothing more nor less than the \enquote{\enquote*{grand, deeply culturally embedded, views of history}} that form the \enquote{\enquote*{trans-historical narrative that is deeply embedded in a particular culture.}}\footcite[][15. Citations omitted.]{Coticchia-2022} Moreover, to borrow the words of a prestigious Israel hater who also happened to have been born an Israeli, it is the \enquote{navigation} of such values, and the \enquote{commitment to liberation,} and the \enquote{particular mode of expression} used in articulating such things that lies \enquote{at the heart of the \enquote*{Theory of Palestine.}}\footnote{Ilan Pappe, \enquote{The Agency and Resilience of the Palestinians Shines: On Imagining Palestine: Cultures of Exile and National Identity by Tahrir Hamdi,} (n.d.), Janus Unbound, \url{https://www.janusunbound.com/book-reviews/imaginingpalestine-jwfka-8hh6a}.}

Almost invariably emphasizing that Palestine's victims are women, children, journalists, doctors or those living otherwise innocent or worthy lives, and making no mention of the deep penetration of Hamas into civil society, civilian infrastructure and below-ground tunnels for combat operational purposes, this epic story of plight is packaged as the \enquote{comprehensive annihilation of a people's social, cultural, intellectual and biological continuity.}\footnote{Ghada Ageel, \enquote{Israel is killing doctors so Gaza can never heal from genocide,} Al Jazeera, July 9, 2025. \url{https://www.aljazeera.com/opinions/2025/7/9/israel-is-killing-doctors-so-gaza-can-never-heal-from-genocide}.} In its totality, as well as in its essence, it constitutes a mode of communication that bundles various manipulative techniques, all designed to invoke feelings of group-specific guilt and having the effect of dehumanizing the group and its allies, which must all be held collectively responsible. Thus, Israel, the United States and \enquote{The West} all stand accused. But while such frenzied narratives may have been for the most part dismissed throughout the second half of the prior century by Western audiences, as contrary to both fundamental honesty and common sense, their deeply troubling ascendance in the current era lies at the heart of this paper's overall thesis. The key point here, perhaps most tellingly, is that this phenomenon is not only the result of the long chain of events and trends described in preceding sections, with the UN at the center of gravity of its continual uplift; rather, it also represents the start of a new era of calumny aimed primarily at Jewish people, and all other groups that stand in the way of what author Bruce Bawer detected years ago and named \textit{The Victims' Revolution.} As a reviewer of Bawer's book noted, it \enquote{confirms what many have long known or suspected: most of what happens in \enquote{identity studies} is pseudo-intellectual because its goals are nonintellectual.}\footnote{Adam Kissel, \enquote{Victims as Victimizers} (review of The Victims' Revolution: The Rise of Identity Studies and the Closing of the Liberal Mind, by Bruce Bawer), Academic Questions 26, no. 1 (2013): 106--110, 106. \url{https://doi.org/10.1007/s12129-013-9333-z}.} 

This section thus turns to how the grotesque imagery and rhetoric now routinely used to define Israel have become the vortex of the victims' revolution. Befitting Israel's purported status alongside the worst regimes of the last century (South Africa's apartheid regime and Germany's Nazi regime), Israel hatred has emerged into a sort of bewildering heyday. It has now become so wildly popular as the sum and substance of the victims revolutionary portfolio that is applied mechanistically, with no real thematic variation, so that the immediately prior century's worst nightmares must be seen as irrevocably attached to a geographically miniscule nation that possesses both a thriving, diverse culture and increasingly prosperous economy. Not only has this master narrative turned the Jews (now updated to the \enquote{Zionists}) from the archetypal victims of the last century into the quintessential villains of the present one, it has embodied an outlook in the present era that reflects a monumental breakdown in reasoning which threatens to engulf the views of vast numbers of well-intentioned and often well-educated people in Western society, such as college professors, non-profit managers and countless others. But for their desire for minor, personalized celebrity in their preferred channels of social media to proclaim that they are on the side of Good, they might have remained anonymous. Yet they are now highly visible, railing on against what they perceive as the Evil.

In this resulting worldview of a harmonious and peaceful world despoiled only by Israel, a seemingly eternal (if occasionally dormant) hostility has thus been resurrected toward Jewish identity, transformed from the prior century in the form of Jews as purveyors of a perverse religiosity\footcite[4 (``In Christian Europe, Jews and Judaism had been defined in terms of religion, not ethnicity, habitus, or civilization. But even as a (superseded) religion Judaism was having a hard time in nineteenth-century Germany. Its collective legalism was framed as the antithesis of personal religiosity (\textit{Glauben} and \textit{Innerlichkeit}), its covenantal tribalism as incompatible with the principles ofHumanität. And whereas Protestant theologians could muster a certain regard forthe pristine Hebraismus of the biblical prophets, they had no sympathy whatsoever for `degenerate' post-exilic \textit{Judenthum}.)'']{Zwiep-2022} into their contemporary characterization as purveyors of a perverse \enquote{ethno-supremacist national ideology.}\footnote{Omar Suleiman, \enquote{The Nakba Explained: Why Israel Won't Win,} May 20, 2024. Yaqeen Institute for Islamic Research. \url{https://yaqeeninstitute.org/watch/series/explained/the-nakba-explained-why-israel-wont-win}. (\enquote{On the anniversary of the first Nakba, we remember the lives destroyed by the ethno-state of Israel, by the dangerous ethno-supremacist ideology of Zionism, and by the continued complicity of the United States government at every single level, despite you not consenting to your tax dollars being used to fund this atrocity. This didn't start on October 7th, this started over 70 years ago.})} Such a characterization destroys any sympathy, legitimacy or goodwill that may otherwise be felt toward Israel and its supporters. Thus, Israel's national future is to be sacrificed to save its Palestinian victims from their otherwise endless \enquote{Zionist problem.}\footnote{G.W.O.O. Palestine and the Zionist Problem (Book Review), July 25, 1920. The New York Times Book Review and Magazine. \url{https://www.nytimes.com/1920/07/25/archives/palestine-and-the-zionist-problem.html}. (\enquote{Palestine is today a land filled with sacred association for the followers of three great religions, and cannot be said to belong to any particular group.})} The growing intensity of this desire for a \enquote{shaking off} of Israel, its Jewish national identity and its supporters ---  which reflects the actual meaning of the Arabic word \textit{intifada}\footnote{Natakallam Blog, \enquote{The Origins and Meaning of Intifada}, (n.d.)  \textsc{url: } \url{https://natakallam.com/blog/meaning-of-intifada/}(``The word comes from the Arabic word 'nafaDa' which is a verb meaning 'to shake the dust off something.' It could be used literally as ``shaking [the dust] off the carpets'' or figuratively, as 'shaking off one's laziness, or ``being finished or rid of something.''' (Arabic omitted.)} --- overpowers all countervailing logic, in the name of reversing what is now universally acclaimed to be  the paragon example of societal and racial injustice. With the mantras of \enquote{anti-Zionism} (but not antisemitism) having thus taken center stage among the world's multitude of \enquote{anti-racist} causes, the question becomes how this absolutist worldview came to capture so much prominence on the dashboard of Western priorities, so that they now resonate as deeply as appeals against South African brutalities or Nazi atrocities only a did several decades ago.\footnote{For an example of an intellectualized comparison between Nazi Germany and Zionist Israel, see Masha Gessen, \enquote{In the Shadow of the Holocaust,} December 9, 2023. The New Yorker. \url{https://www.newyorker.com/news/the-weekend-essay/in-the-shadow-of-the-holocaust}.}

The answer may be akin to a phrase in the Ernest Hemingway novel \textit{The Sun Also Rises.} When asked how he went bankrupt, a character replies: \enquote{Two ways. Gradually, then suddenly.}\footnote{Cardinal Team, August 12, 2021. Cardinal Institute for West Virginia Policy. \url{https://cardinalinstitute.com/gradually-then-suddenly/}.} As noted in the preceding section, after the Oslo peace process stalled in the late 1990s, efforts aimed at Israel's delegitimization were gradually ampilifed by civil society actors, together with states that continued to harbor hostility toward Israel, thus resuming the Resolution 3379 era of accusations of racism that were ostensibly put to rest in 1991 with UNGA resolution 46/86.\footnote{Elimination of racism and racial discrimination : resolution / adopted by the General Assembly. December 6, 1991. \url{https://digitallibrary.un.org/record/135193?v=pdf}} The process began in earnest at the World Conference Against Racism in 2001 (Durban I), continued through the BDS campaign's launch in 2005 by NGOs complicit in the resistance movement,\footnote{  Palstinian Civil Society Organisations. \enquote{Palestinian Civil Society Call for BDS.} BDS Movement, August 7, 2005. \url{https://bdsmovement.net/call}.} and received a bold endorsement in the opening speech of Iran's Mahmoud Ahmadinejad at the Durban II conference in Geneva, in 2009.\footnote{al Majdal Editorial, Autumn 2011,\enquote{Israel and the Crime of Apartheid: Toward a Comprehensive Analysis, }\textit{BADIL Resource Center for Palestinian Residency and Refugee Rights}. \url{https://badil.org/phocadownload/Badil_docs/publications/al-majdal47.pdf}, pp. 2-3. (\enquote{In the face of the peace process's inability to realize Palestinian rights, civil society actors and activists sought an analysis that could encompass and describe the full range of Palestinian grievances and oppression. This effort crystallized, most notably, during the international civil society forum that met in Durban in 2001 in a conference running parallel to the official UN World Conference Against Racism.})} The efforts continued into the mid- to late-2010s, with a formal call to \enquote{globalize the Intifada} in 2014,\footnote{Shideler, K. and Daoud, D. (8/19/2014). \textit{Command and Control: The International Union of Muslim Scholars, The Muslim Brotherhood, and The Call for Global Intifada during Operation Protective Edge}. Center for Security Policy Occasional Paper Series. \textsc{url: } \url{https://www.centerforsecuritypolicy.org/wp-content/uploads/2014/08/Command_and_Control_08-21-14_12.50pm.pdf}.} the so-called \enquote{Third Intifada} of 2015--16,\footnote{Yumna Patel, \enquote{A decade in review: the moments that shaped the past 10 years in Israel/Palestine,} December 24, 2019, \url{https://mondoweiss.net/2019/12/a-decade-in-review-the-moments-that-shaped-the-past-10-years-in-israel-palestine/}.} Hamas' \enquote{Great March of Return} in 2018-19.\footnote{Jane Ferguson, ``Gazans suffer life-shattering injuries when border protests turn violent,'' PBS News Hour, August 20, 2019. \url{https://www.pbs.org/newshour/show/casualties-of-war-2}.} and the May 2021 conflict in Gaza.\footnote{IMEU, Palestine: 2021 Year in Review, \url{https://imeu.org/resources/resources/palestine-2021-in-review/119}.}

Considering the line in Hemingway's novel, the years between September 2001 and October 2023 may thus be considered the \enquote{gradually} phase; the era since then is the \enquote{suddenly} one. But by parsing the chronology of the resurgence of Israel hatred in this fashion, a deeper issue arises: what laid the foundation for it in the first place? That is, how was it that long after 1991, when the \enquote{Zionism is racism} formula lost its formal backing in the UNGA, the seeds for its return were hardy enough to take root again and return in a new era of what one sympathetic commentator has described as \enquote{revolutionary mobilization}?\footnote{Zack Cuyler, \enquote{Exploitation, Elimination, and Rebellion in Israel/Palestine: Introduction,} International Labor and Working-Class History. 2025;108:326-328, 327. \url{https://doi.org/10.1017/S0147547925100082}.} By way of comparison, during the era leading up to and including the Oslo process, there were only faint rumblings that Israel should be denounced as an apartheid state. One example is found in a 1983 pamphlet in which Israel's reliance on captive Palestinian labor was noted as a fundamental characteristic of what can be considered (as noted in section II) an intrinsic element of apartheid as an archetypal system.\footnote{Elisabeth Mathiot, \enquote{Zionism, A System of Apartheid}, 1982. EAFORD Paper No. 25, 4. \url{https://www.eaford.org/site/assets/files/1023/zionism-_a_system_of_apartheid.pdf}.} Published by EAFORD, an NGO deeply engaged in anti-Israel advocacy at the UN which had ties to Moammar Kaddafi's regime in Libya,\footnote{Anne Bayefsky, Sponsor of UN-hosted blood libel tied to Libyan regime, April 18. 2010. UN Watch. \url{https://unwatch.org/sponsor-of-un-hosted-blood-libel-tied-to-libyan-regime/}.} its author had earlier taken part in France's anti-apartheid movement after the Algerian war of independence.\footnote{Tal Sela, \enquote{Not Merely a Newsworthy Commodity: Jean-Paul Sartre's Engagement in the Struggle Against Apartheid,} Critical Arts, 34(1), 41--55 (2020). Notes 12 and 13. \url{https://doi.org/10.1080/02560046.2019.1691616}.} But by the early 2010s, at a time when Israel and Palestine had long been considered on separate national trajectories as a result of the Oslo process (unlike the situation in South Africa, where Bantustan political independence was never seriously contemplated), Israel's oppression was alleged by John Dugard --- a South African academic who was directly involved in UNGA committee work relating to the conflict --- as \enquote{even worse} than South Africa's because Israel did less for Palestinians than South Africa did for those it consigned to Bantustans.\footnote{John Dugard, The Law and Practice of Apartheid in South Africa and Palestine. Al-Majdal Magazine, 47 (2011)
\url{https://badil.org/publications/al-majdal/issues/items/1411.html}.}  To understand how an Israel-South Africa equivalency had thus moved from a sideshow in the polemics of an NGO on the UN's fringes (that is, EAFORD) to become part a primary focus of the UNGA, warranting its passage of the \enquote{Zionism is racism} resolution,  two aspects of the transition require further examination. The first is the final years of the career of Dr. Fayez Sayegh (1922-80), a pan-Arab diplomat and author whose hatred toward what he once described as \enquote{the Zionist program of incursion into the Arab World}\footnote{Fayez Sayegh, \enquote{Neutralism in the United Arab Republic}, in: The Dynamics of Neutralism in the Arab World (F. Sayegh, ed.), (Chandler Publishing Company San Francisco (1964)) 163--224, 169.} provided the impetus for his claims of racism toward Palestinians.\footnote{The second aspect --- which addresses how the UN's direction on dealing with racism was heavily influenced by both incessant anti-Israel advocacy and the rise of critical, \enquote{anti-racist} inquiry in the Western academy --- will be delved into the next, and final, section of this paper.} 

As noted in section I, Dr. Sayegh's 1965 monograph, \textit{Zionist Colonialism in Palestine,} is the foundational document of the \enquote{Zionism is racism} equation, and it has organized political support even today, notwithstanding the withdrawal of Resolution 3379 nearly 35 years ago.\footnote{Sabrina Miller, \enquote{Greens Fight over Whether \enquote*{Zionism Is Racism.}} The Telegraph, June 6, 2026. \url{https://www.telegraph.co.uk/politics/2026/06/06/greens-fight-over-whether-zionism-is-racism/}.} The monograph was published shortly after the PLO launched the ``Palestine Research Center,'' which it placed under Sayegh's direct editorial oversight, even though he apparently remained employed in a diplomatic capacity for Arab states at the UN. In it, Sayegh depicted Israel as not only an unwelcome interloper, but also an inherently racist one. Why did he shift his critique of Israel to center on racism after having written extensively about it in prior articles without giving any attention to this particular societal condition? It may have been that, like other Arab intellectuals living abroad during this era, he saw himself as holding \enquote{an interstitial position between two worlds,} giving him an opportunity to be among those who were \enquote{catalysts for change.}\footnote{Zeina G. Halabi, \textit{The Unmaking of the Arab Intellectual,} Edinburgh University Press, 20170, 103.  At the time he edited \textit{The Dynamics of Neutralism in the Arab World} in the early 1960s, Sayegh was at St. Antony's College, Oxford University, Oxford, England, on a research grant.} He may have been influenced by Black-Algerian Marxist theorist Frantz Fanon, whose highly influential 1963 book, \textit{The Wretched of the Earth}, generated Arab interest in the early 1960s, soon after publication.\footnote{
Mounir Saidani, \enquote{On Reclaiming Fanon: From and for the Arab Maghreb}, In: Colonial Legacies and Arab-Majority Regions, 141--59. Bristol University Press, 2026. \url{https://doi.org/10.51952/9781529240566.ch006}.} Had he investigated further, he would have easily found \enquote{Racism and Culture,} a speech Fanon gave at the first Congress of Negro Writers and Artists in Paris in 1956.\footnote{Frantz Fanon, \textit{Toward the African Revolution: Political Essays}. (Grove Press, 1967, Trans, Haakon Chevalier), 31.} In it, Fanon established the equation: \enquote{In reality, a colonial country is a racist country.}\footnote{Id. at 39--40.} He might also have pondered Fanon's point that \enquote{At the United Nations there is a commission to fight race prejudice.}\footnote{Id. at 39.} Or, he may have recalled how the Soviet delegation challenged a proposed American insertion into the draft ICERD with countervailing language that \enquote{policy and ideology of colonialism, national and race hatred and exclusiveness,} as discussed in section II.

In any case, as his career progressed, and he became a fixture in the Arab delegation offices at the UN, Sayegh continually refined his monograph's claim that racism was at the heart of what it described as \enquote{the Zionist settler state.} His emphasis on Israel's original sins of colonialism and racism was also repeated almost word-for-word in a statement echoing a Fanonian universal argument that was ``presented by the General Union of Palestine Students.'' Entitled ``The Palestine Problem.''\footnote{Palestine Problem: Statement presented by The General Union of Palestine Students, \url{https://collections.lib.utah.edu/details?id=841869}. The digital catalog describes this document as a ``Statement equating Zionism with racism and listing violations of Palestinian rights; undated but probably from ca. 1965-1967.''} Issued not long after his 1965 monograph was published, it appears likely to have been the product of his ghostwriting pen:

\begin{adjustwidth}{.65cm}{.5cm}
  \begin{singlespace} 
    [A]s a racist system animated by doctrines of racial self-segregation, racial exclusiveness and racial supremacy, and methodically translating these doctrines into ruthless practices of racial discrimination and oppression — the political systems erected by Zionist colonists in Palestine cannot fail to be recognized as a menace by all civilized men dedicated to the safeguarding and enhancement of the dignity of man. \textbf{For whenever and wherever the dignity of but one single human being is violated, in pursuance of the creed of racism, a heinous sin is committed against the dignity of all men everywhere.} (Emphasis in original.)
  \end{singlespace}
\end{adjustwidth}

By this time, Sayegh's focus had clearly shifted. No longer was he content, as he had been in the 1950s, to critique the left-leaning nature of Israeli society \parencite{Sayegh-1958b} or glorify the unalterable nature of Arab unity in the Levantine world of Syria and Lebanon (which included Palestine as \enquote{southern Syria} due to \enquote{the occupation thereof by Muslim-Arabs from the Arabian Peninsula in the Seventh Century AD} \parencite[12]{Sayegh-1958a}. As Sayegh took his place as an eminent figure in the rarefied world of international diplomacy of all Arab states --- a far cry from his early years as a political operator on the fringes of a chauvinistic Syrian movement that was ultimately suppressed by the Assad regime\footnote{Hussein Aboubakr Mansour, ``The Man Who Made Zionism Into Settler Colonialism: Antoun Sa'adeh, Fayez Sayegh, and the Intellectual Genealogy of Anti-Zionism,'' November 13, 2025, \url{https://www.meforum.org/mef-online/the-man-who-made-zionism-into-settler-colonialism}.} --- he was now directing all his energy toward championing the cause of Palestine on behalf of the Arab bloc of UNGA members states.

Over the final decade of Sayegh's life, while serving ``in various capacities for a variety of Arab states at the UN, and ultimately as chief of the Arab States' delegation to the UN,''\footnote{E.g., Abrahamian 2009.} his contribution to normalizing Israel's worldwide ostracism reached its climax. He used his deep experience in navigating the UN's many forums, systems and committees to justify the Arab states' outright refusal to engage directly with Israel, because it was entirely up to the UN to solve the \enquote{whole problem} it had created by recognizing Israel in the first place.\footnote{Kadhim Jawad, ``Selective Debates on Palestine,'' Baghdad Magazine, 1970, \url{https://archive.org/details/selective-debates-on-palestine/}, p. 11.} Meanwhile, to make the UN was as hostile as possible to Israel, so that her ability to engage diplomatically in its forums was poisoned at the base, Sayegh became ``the principal author'' of Resolution 3379\footnote{The Washington Post (Dec. 14, 1980). Fayez Sayegh, Leader of Arab Groups In U.S., U.N., Dies in New York City. (Online; accessed Aug. 8, 2023.) \textsc{url: } \url{https://www.washingtonpost.com/archive/local/1980/12/14/fayez-sayegh-leader-of-arab-groups-in-us-un-dies-in-new-york-city/c3640054-e9d5-4cc4-9f96-525e350d4df1/}} while serving as ``Senior Consultant to the Kuwait Ministry of Foreign Affairs and as a member of the Delegation of Kuwait to the United Nations.''\footnote{Fayez A. Sayegh collection on the Arab-Israeli conflict:  Overview of the Collection (n.d.). Archives West. \url{https://archiveswest.orbiscascade.org/ark:80444/xv382707}.} Throughout this time, he also traveled frequently, giving numerous lectures and presentations to groups interested in learning about the Arab world's views on Israel and other topics. He also represented Kuwait on the UNGA committee established under Article 8 of the ICERD (that is, the CERD) and appears to have served as its ``rapporteur'' from 1968 until his death in 1980.\footcite[][Sayegh's rapporteur role in CERD  from 1968 until his death in 1970 is also noted without reference to any source in Feldman 2015, 37.]{Sayegh-1982}. The CERD was \enquote{the first body established by the United Nations to examine the application of a human rights treaty.}\footnote{The Practical Guide to Humanitarian Law (n.d.) Doctors without Borders, \url{https://guide-humanitarian-law.org/content/article/3/committee-on-the-elimination-of-racial-discrimination}/.} And, as is shown by the role of EAFORD, which \enquote{derive[d] both its name and its inspiration} from the ICERD, it was always the goal of Arab states to use the ICERD in a way that effectively weaponized its overall purpose against Israel.

The racism theme thus established by Sayegh's work --- both as a writer for the PLO and as the author of the UNGA's Resolution 3379 --- was also not confined to the official deliberations of general assembly or the various committees and task forces of the UN. It was also socialized  within the UNGA by Arab states, as shown by a letter from Iraq's UN ambassador which distributed the resolutions of a ``Symposium on Racism'' held the following year in Baghdad, as an annex to a UNGA agenda item.\footnote{Comprehensive Review of the Whole Question of Peace-keeping Operations in all Their Aspects --- Elimination of all Forms of Racial Discrimination, Annex to A/31/339, November 22, 1976), \url{https://digitallibrary.un.org/record/153923/files/A_31_339-EN.pdf}.} The symposum's resolutions noted that since Israel's creation in 1947 on the basis of a proposed partition of Palestine, ``the resolution of November 1975 ... had been adopted when membership of the United Nations had become more genuinely representative of the opinion of the world as a whole''; that ``Zionism as a colonial settler concept was an offshoot of 19th century imperialism''; and that ``[n]othing is more dishonest than the slogan --— unleashed by the United States and Israel as the principal weapon in their campaign against the decision of the United Nations --— that anti-Zionism is anti-Semitism.'' But perhaps the most incisive argument brought forward by Israel's enemies to delegitimize Israel can be found in the declaration issued by the forum held in Tripoli, Libya in the spring of 1976, at which Sayegh spoke and probably helped edit, if not write:\footnote{Editor's Preface to Fayez Sayegh, \textit{Racism and Racial Discrimination Defined}, 1982.}

\begin{adjustwidth}{.65cm}{.5cm}
  \begin{singlespace} 
  The ruthless campaign of vilification to which the United Nations has been subjected, by Zionism and its racist and imperialist allies, in response to its principled and courageous adoption of the afore-mentioned resolution reveals the desperation and isolation of the forces of racism and Zionism in today's world. Nothing is more dishonest than the slogan —    unleashed by the United States and Israel as the principal weapon in their campaign against the decision of the United Nations — that anti-Zionism is anti-Semitism.
  \end{singlespace}
\end{adjustwidth}

\noindent
The fact that the written conclusions of both the Tripoli and Baghdad conferences emphasized that \enquote{anti-Zionism} could not possibly be considered anti-Semitism is also notable here. It directly echoes the assertion made by the Soviet delegation in its proposed amendment to the wording of ICERD, reviewed in section 2, that sought to differentiate between anti-Semitism --- which ended up \textit{not} being specifically mentioned as a form of racism in the wording of the ICERD --- and Zionism, which could be equated with ``colonialism, Nazism and neo-Nazism.''\footnote{Section 2, fn. 46.} In effect, Resolution 3379, which operationalized Sayegh's 1965 monograph, reverted to the Soviet interpretation of the ICERD as reflected by its proposed amendment to its wording, even though that amendment was not in its final language. Thus, by the mid-1970s, a premeditated attempt by Arab state delegations within the UN --- Libya, Kuwait, Iraq, and others, in close coordination with the PLO and in conjunction with eastern bloc and African states --- was in full swing. Its goal was to depict Israel and its western allies as flagrantly racist colonialists who exploited subjugated peoples.\footnote{See, e.g., United Nations, General Assembly 13th session, Special Political Committee, Summary Records, 978th meeting. New York NY: UN Headquarters, November 17, 1975. Agenda Item 54 - United Nations Relief and Works Agency for Palestine Refugees in the Near East (continued) (A/10114, A/10115, A/10268). UN Document A/SPC/SR.978, p. XXX. \url{https://digitallibrary.un.org/record/1307549}. Remarks by Mr. SIBAHI (Syrian Arab Republic). (\enquote{The Committee was currently being asked to deal with a situation for which the racist, Zionist entity and its allies and supporters, notably the United States of America and certain Western States, were fully responsible.}} In the long run, it is now clear that this strategy was meant to establish that just as the racist practices of South Africa were abhorrent to the world, so too should Israel's be --- and the West's. After all, the special value of connecting Israel (always identified as ``Zionism'') with ``other racist regimes, as evidenced by its close relationships with Rhodesia and South Africa, [was in reinforcing it as] \dots a natural outcome of its roots and development, for it has always drawn its support and sustenance from imperialism and settler-colonial regimes.''\footnote{Palestine Information Office, International Declarations --- Zionism and Racism. 1981, p.15. This booklet contained \enquote{declarations and resolutions \dots adopted during two international symposiums held in 1976 in Tripoli, Libra and Baghdad, Iraq.}} 

But notwithstanding this manipulation of the purpose of a UNGA convention on racial discrimination to direct its most far-reaching inquiries and accusations against Israel, on a basis that eventually made it equivalent to apartheid-era South Africa, it is not altogether clear that the reality of settler-colonialist oppression was what truly what drove Sayegh and many others in the Arab world to denounce Israel so categorically. It may have been useful to frame for both Arab and non-Arab audiences the basis of reasoning that Westerners could recognize was derived from their own disgraceful histories. But a more fundamental reason for Sayegh's enmity toward Israel seems to have centered on the resurgence of a Jewish national identity that it demonstrated, which stood apart from and beyond Jewish religiosity. In so doing, it could no longer be consisdered secondary to the supremacy of Islamic religiosity and national identity in the MENA region, reflected by its \textit{dhimmi} traditions. This is evident from at least two of lines of argument used by Sayegh, EAFORD and others in the Arab world's  anti-Israel camp. The first can be gleaned from Sayegh's early writing on Arab unity, as expressed in his 1958 book, \textit{Arab Unity: Hope and Fulfillment,} in which the erasure of any Jewish role in the history of the MENA region is described with a sense of ultimate finality in a grand accomplishment, notwithstanding the eternal vicissitudes of human history. In describing the process of conquest of geographies adjacent to the Arab peninsula, he wrote:\footcite[16]{Sayegh-1958a}

\begin{adjustwidth}{.65cm}{.5cm}
    \begin{singlespace}
    [T]he stamp which was in due course imprinted on the entire population, and which was destined to prove permanent, persisting until today, came from the Muslim-Arab invaders. For, while ethnically and culturally they underwent mutual absorption with the compound-population they conquered, they nevertheless bequeathed their language and their new faith to the amalgam. Not only the Arabs themselves, then, but all their predecessors as well, who were arabized by them, `were and remain.'
    \end{singlespace}
\end{adjustwidth}

Apart from his obviously deep sense of pride in his Arab heritage, Sayegh --- who was the son of a Presbyterian minister\footnote{Andrew I. Kilgore, \enquote{Years After His Death, Dr. Fayez Sayegh's Towering Legacy Lives On,} Washington Report Archives (2000-2005), December 25, 2005. \url{https://www.wrmea.org/2005-december/in-memoriam-25-years-after-his-death-dr.-fayez-sayeghs-towering-legacy-lives-on.html}.
} --- could also not suppress his natural antipathy toward anything that might tend to compete with collective glory of he Arab world, including Israel. This is apparent from his unsuppressed desire to mock Israel because its Jewish citizens expressed admiration of their own religious concept of \enquote{chosenness.} In his remarks at the Tripoli symposium in 1976, Sayegh elaborated on his disdain for Zionism by dwelling on the reliance of some of its adherents to this concept, about which he said \enquote{zionist literature abounds ... [reflecting] this sense of superiority by the zionists.}\footcite[5-6.]{Sayegh-1982} Actually, the concept abounds in Judaism based on ancient holy texts, where the Jews considered themselves having the promise of ultimate redemption notwithstanding their many sins by their God's gracious choice. Moreover, in and of itself, the sense of pride in Judaism reflected by being \enquote{chosen} says nothing about the likely treatment of Arab inhabitants. Numerous multicultural societies have population segments that may seem to take excessive pride of their own ethnicity, whithout that translating into legalized discrimination against members of other cultures on racist grounds. Such a charge of state racism was not even a widely held view among Arab scholars prior to Sayegh's monograph. Even Sayegh, in an earlier time frame, had not chosen to describe Israelis as Westernized colonialists but rather as purveyors of the \enquote{Communism [which] exercises [its sway] over the minds, hearts and aspirations of large sectors of the Israeli people} \parencite{Sayegh-1958b}.

In fact, well into the 1970s, Palestinian scholars never made much of a point of chronicling a long history of Zionist racism. For example, in a 1977 article entitled ``The Palestinian Position,'' the well-known historian Sabri Jiryis (who began a long tenure as director of the Palestine Information Center the following year) chose to describe \enquote{racialist Zionist principles} by referring to Israel's refusal to withdraw from areas territories taken during the 1967 war in order \enquote{to keep them open to Jewish settlement, in addition to making them a market for Israeli goods and a reservoir of labour forces that may be required by her economic development.}\footnote{Sabri Jiryis, The Palestinian Position (1977). Journal of Palestine Studies, Vol. 6, No. 4, 150--158, 153. \url{https://www.jstor.org/stable/2535787}. The point itself again highlights how control over external labor was essential to a \enquote{racialist,} and ultimately apartheid, analysis.} This could also be explained as ensuring Israel's continuous ability to maintain a Jewish presence on long that lies at the heart of ancient Judaism. But also in the late 1970s, when Jiryis published Volume I of his magnum opus, the \textit{History of Zionism,} (covering the period from 1862 to 1923), no real mention is made of racism. It was not until Volume II (published in 1986) that it a brief reference appears in noting that the leaders of the Second Aliyah of the 1920s sought to hire Jewish labor to solidify the communal growth of the \textit{Yishuv}.\footnote{Sabri Jiryis and Fida Jiryis, \textit{The Foundations of Zionism}, Ebb Books (2025). Ebb Books, which describes itself as \enquote{an independent press that critically assesses British politics from a communist perspective.} (\url{https://www.ebb-books.com/about} chose to publish this book on October 7, 2025.} When it was finally published in a translated English edition in 2025 and renamed \textit{The Foundations of Zionism,} the more precise concern with societal racism appears only in a ``newly written conclusion'' entitled \enquote{Zionism in the service of colonialism}\footnote{Ebb Magazine. \enquote{The Foundations of Zionism — Ebb Magazine.} Accessed July 6, 2026. \url{https://www.ebb-magazine.com/books/p/the-foundations-of-zionism}.} In it, the now 88-year-old author describe how he was treated like a ``second-class citizen'' even though he was born in 1938, had a residence in the same area where he was born for over at least a decade after 1948, graduated Hebrew University in the late 1950s, and practiced law in Haifa before concerns over his political activities led to his expulsion in 1970. The conclusion further notes that \enquote{Zionist racism never disappears, it merely changes form,} because \enquote{those racist attitudes towards the Palestinians in Israel are only an extension of the Zionist position towards the Palestinian people as a whole, as the settlers attempt to fragment it and replace it with themselves.}\footnote{Id., Conclusion.} If racist \enquote{attitudes} and broad claims of \enquote{displacement} are sufficient to justify an excommunication of Israel from the world's family of naitons, surely there are other nations that qualify to offer it company.

Before reaching conclusions about Mr. Jiryis' apparent conversion to the view that Zionism and Israel are founded on ``racist attitudes,'' it is also worth pausing to consider the story of his family after it left Israel in the 1960s so that he could take a senior role in the development of PLO media communications. He was out of the country until the consummation of the Oslo Accords, after which he was permitted to return to his native village in the Galilee, where he has kept a residence ever since.\footnote{Peter Beinart, ``We are not wanted there,'' February 22, 2026, The Beinart Notebook. \url{https://peterbeinart.substack.com/p/we-are-not-wanted-there}.} While he may continue to believe that Israeli Arabs are disenfranchised, and his daughter and English translator Fida --- who also returned to ``historic Palestine'' --- apparently agrees,\footnote{Fida Jiryis: A Unique Story of Return, 2017, Al Jazeera, \url{https://interactive.aljazeera.com/aje/2017/palestine-in-motion/fida.html}. (``We can vote; we have access to education, employment, healthcare and social benefits; we are free from Israeli occupation, military checkpoints, army incursions and Jewish settler violence endured by our brethren in the West Bank and Gaza. Yet, underneath this facade is a complex system of discrimination in every aspect of our lives.'')} his personal story actually reflects a far more nuanced than one of unmitigated racist oppression. It is also consistent with the experiences of many present-day Arab Israelis, such as Nuseir Yassin, an Arab Israeli entrepreneur and media personality whose family remained in Israel after 1948. In summary on this point, True, there are societal problems that remain a challenge for Israel's Arab sector, which makes up roughly 20\% of the entire national population, but they can hardly be considered the type of discrimination that would otherwise lead to the inexorable conclusion that Israel must be treated as the lone pariah among the world's nations. While there are numerous distinctions in the Israeli system based on what is described as a citizen's original ``nationality'' --- whether it be Druze, Arab, or Jewish --- and notable disparities in living standards among them, the overall reality reflects an imperfect system in which alleviating hardship associated with the adversities borne by ethnic minorities is considered a widely-shared societal concern, much as is true in many other advanced nations. And yet unlike those countries, there is nonetheless an ongoing uproar over Israel's racially discriminatory shortcomings. With similar problems surely evident elsewhere, the question that must be asked is why is this now singled out so dramatically in relation to Israel and virtually no where else in the MENA region? 

There are a few paragraphs missing here relating to Arab perceptions of Israel as an affront to \enquote{authentic Islam} of the Arab world and its effect on Arab nationalists who may have given up on the pan-Arab dream but remain inspired to view the \enquote{return} of Palestine to the Arab fold as critical to restoring the direction of the Arab world. The \enquote{racism,} \enquote{apartheid}
and \enquote{genocide} angles should ring hollow if the concept of what really drives the continuous preoccupation with Israel among Islamist intellectuals, but they have succeeded in repackaging it as a crime perpetrated on the West for these reasons, even thought their real passion for the cause is derived from how it breaks the concept of the Arab world as a unified, indivisible whole, which is an exemplar to the rest of the Islamic world. Sources for this are found in Zaina Halabi's 2017 book, \textit{The Unmaking of the Arab Intellectual}\footnote{Zeina G. Halabi, \textit{The Unmaking of the Arab Intelectual: Prophecy, Exile, and the Nation,} Edinburgh University Press, 2017} and Ibrahim Abu-Rabi's \textit{Contemporary Arab Thought: Studies in Post-1967 Arab Intellectual History.} \parencite{Abu-Rabi-2004}

Resuming regular programming

In seeking conclusions on the accelerating rise of Israel hatred in today's world, it seems clear that goading the West into casting judgment of eliminationist guilt on the Jewish state on the grounds of its inherent racism has become an overarching objective. Noura Erekat has all but described it as the key ingredient of a ``dominant Palestinian tradition,'' derived from ``Palestinian intellectual thought,'' that Zionism as \enquote{a form of racism and racial discrimination} must be established and recognized as a foundational and unassailable truth, regardless of whether the UNGA backtracked on its resolution to the same effect.\footnote{Erekat, N. (5 July 2021). Beyond Discrimination: Apartheid is a Colonial Project and Zionism is a form of Racism. \textit{EJIL: Talk! Blog of the European Journal of International Law.} \url{https://www.ejiltalk.org/beyond-discrimination-apartheid-is-a-colonial-project-and-zionism-is-a-form-of-racism/}.} Borrowing a well-known phrase coined by Black American socialist (and eventual Marxist) W.E.B. Du Bois,\footnote{Jeff Goodwin, \enquote{Socialism was central to W. E. B. Du bois's thought,} July 2, 2026. Jacobin. \url{https://jacobin.com/2026/07/du-bois-black-marxist-tradition}.} this tight coupling of Israel and Racism depends on ``analyzing the workings of the color line in Palestine and Israel and mapping the trajectories through which race and racialization came to play a central role in the land, illuminate the truly global reach of its logic and highlight some of the properties that have made the global color line so durable.''\footnote{A Bivins, Tracing the Historical Relevance of Race in Palestine and Israel, August 4, 2021. \textit{MERIP}.  \url{https://merip.org/2021/08/tracing-the-historical-relevance-of-race-in-palestine-and-israel/}.} 

But while the ``color line'' concept can easily be seen as intellectually honest and potentially useful to describe 20th century Black experiences in the Jim Crow southern United States, and in the totalitarian effects of apartheid South Africa, its application to Israel should remain more complicated than cutting and pasting those historical realities on today's Palestinians. As explored in previous sections, the fundamental problem is not just whether they can be deemed \enquote{racialized} based on something other than skin color, or on what their own champions describe as \enquote{traditional conceptions of \enquote*{race}} \parencite[885]{Dugard-2013}. but whether the proscriptions of international law would extend to such a form of discrimination at all. As at least one more honest scholar in the same ideological mold has recognized, if the treaties and conventions applicable to apartheid and genocide are to be respected in the era in which they were passed, being \enquote{Palestinian} cannot be the basis of racialization because the definitional framworks were \enquote{grounded in the assumption that \enquote*{race} is immutable and biologically determined.} \footnote{Lisa Davis and Kirby Anwar, \enquote{Race, Gender and Apartheid in the Draft CAH Treaty: South Africa's Role,} OpinioJuris, July 18, 2025, \url{https://opiniojuris.org/2025/07/18/race-gender-and-apartheid-in-the-draft-cah-treaty-south-africas-role/}} In other words, Palestinian cannot mean both a racial group and a subset of Arab ethnic identity unless \enquote{how international criminal law defines protected \enquote*{groups}} is made clearer.

The fundamental question again arises, which will be turned to in this paper's next and final section: why has racism become virtually synonymous with discussions about Zionism and Israel, and what role has the UN played in this current reality? Clearly, Frantz Fanon was right in noting that \enquote{[a]t the United Nations there is a commission to fight race prejudice,} but how is that among all the countries in the world, Israel is the one constantly being singled out as the incorrigible pariah state on that account? The deeper issue here is whether in the context of Palestine's \enquote{perpetual resistance,} in respect of which Hamas and Palestinian Islamic Jihad clearly now lead the political struggle, Israel's racism is what truly supports their vision of a Palestine that aligns with Western notions of liberal democracy, decency, justice and fairness. Or is it far more closely aligned with Islamist notions of sharia theocracy? Whatever the purported failings of Israel, when held up to the standards of the style of governance that formed the paradigm for the basic concept of a UNGA member \enquote{nation-state,} it is hard to reconcile the charge against the reality.\footnote{Eric L. Jones, The European Miracle: Environments, Economies and Geopolitics in the History of Europe and Asia. 2. ed., Repr, Cambridge University Press, 2003.(\enquote{The Nation-State is nowadays the unit of affairs. It is a purely European form which has been exported to parts of the world that had hitherto known only tribalism. We may look on national income accounting as if it deals in categories both natural and proper, but nation-states are not God-given. They are the creations of post-feudal Europe.})} Instead, it must be understood as something diffferent, which is a grounds invented for discarding the overall contours of the concept of nationality altogether, and replacing it with a style that is exactly the opposite to what it claims it is designed to do. The \enquote{tradition} itself, rather than its target of intolerance, is what is hostile to the diverse lifestyles that result in multicultural societies. How is that, only a few years after the seismic shift of October 7, the idea of the destruction of Israel is still celebrated when it should be seen as antithetical to the very concept of discrimination on which it is based?

One final anecdote on this dubious barometer of the twisted direction of social justice and racial injustice is worthwhile to end this section. A few months before October 2023, a member of Congress who is a leader of the \enquote{progressive} caucus among its Democratic Party members apologized for referring to Israel as a ``racist state.''\footnote{M.Pengelly, Progressive US congresswoman apologises for calling Israel `racist state,' July 17, 2023. \textit{The Guardian}. \url{https://www.theguardian.com/us-news/2023/jul/17/pramila-jayapal-israel-racist-state-apology}.} In her defense, she noted that the phrase came to mind onstage as she sought ``to defuse a tense situation ... where fellow members of Congress were being protested.'' But a New York Times columnist ridiculed what she called the \enquote{hysterical overreaction to Jayapal's \enquote*{racist state} gaffe,} noting \enquote{if Israel's champions are truly worried about the fallout from accusations of racism, they might act to make them seem less credible.}\footnote{Michelle, Goldberg, \enquote{The Hysterical Overreaction to Jayapal's \enquote*{Racist State} Gaffe,} The New York Times, July 17, 2023. \url{https://www.nytimes.com/2023/07/17/opinion/israel-racist-state-pramila-jayapal.html}.} Roughly 18 months later, grassroots organizations aligned with the party's \enquote{progressive} wing  --- a wing that lately appears to have subsumed by the Democratic Socialists of America (DSA) --- went on record with the following:\footnote{ North Carolina Triangle DSA (chapter of the Democratic Socialists of America), \enquote{The Importance of Being Anti-Zionist,} NC Triangle DSA, February 12, 2025. \url{https://www.dsanc.org/leftangles/2025/2/12/4fchhy04n204kn7lrp1mu7e3quj9ge}.}

\begin{adjustwidth}{.65cm}{.5cm}
    \begin{singlespace}
    In the DSA, we do not support Zionism. To take an anti-Zionist stance is to speak and act against the racist violence of the ethno-nationalist state, unleashed with brutality and impunity on the Palestinian people, land, culture, and future and propped up by colonial logics and imperialist calculations of the global powers such as the US. 
   \end{singlespace}
\end{adjustwidth}

This statement makes clear that from a larger perspective, the theme now running so deeply about Israel's racism is not entirely about Israel, even if the takers of \enquote{an anti-Zionist stance} consider that particular postage-stamp sized nation their immediate adversary. Rather, the larger target is revealed by the interwoven narratives that proclaim a perpetual crisis of racism throughout the Western world --- the \enquote{colonial logics and imperialist calculations} ---  resulting from a long, disgraceful history. The echoes of this sentiment can be heard from decades ago. In \textit{Orientalism}, Edward Said wrote that under the ``system of truths'' imbued in Western thought, ``every European, in what he could say about the Orient, was consequently a racist, an imperialist, and almost totally ethnocentric'' \parencite[204]{Said-1978a}. The same \enquote{germ of truth,} if it can fairly be called that, has now crossed the Atlantic Ocean and infiltrated the world's largest liberal democracy with calls-to-action based on its presumed self-evident truth. 

But again, to what extent --- and why --- has the UNGA become the willing institutional handmaiden of this expanding state of destructive, dysfunctional, and perhaps ultimately suicidal state of nation-state affairs? And is it possible that rather than \enquote{dehumanizing the Zionist problem,} that body can at some point reclaim the noble vision of the era of its initial formation, and start working toward \enquote{rehumanizing the nation-state problem}? That is the subject of the next and final section of this paper.

\parindent 0pt
\parskip 6pt

\setenotez{
  list-name = {Notes to Chapter \thechapter}
}
% Place this macro at the exact end of every chapter:
{\footnotesize
\printendnotes[paragraph]}

\thispagestyle{empty} % makes a blank page

\end{document}